\documentclass[12pt,a4paper]{article}
\usepackage[margin=25mm, headheight=15pt, headsep=10pt]{geometry}
\usepackage{fontspec}
\usepackage[table]{xcolor} % Note: table option is required for rowcolors
\usepackage{titlesec}
\usepackage{tocloft}
\usepackage{fancyhdr}
\usepackage{booktabs}
\usepackage{tcolorbox}
\tcbuselibrary{skins, breakable}
\usepackage[colorlinks=true, linkcolor=emerald, urlcolor=emerald, bookmarks=true]{hyperref}

\definecolor{emerald}{RGB}{0,102,51}
\definecolor{darkred}{RGB}{139,0,0}
\definecolor{lightgray}{RGB}{245,245,245}

\usepackage[nil]{babel}
\babelprovide[import, main]{bengali}
\babelprovide[import]{arabic}
\babelfont{rm}[Renderer=HarfBuzz,Scale=1.0]{Noto Serif Bengali}
\babelfont{sf}[Renderer=HarfBuzz]{Noto Sans Bengali}
\babelfont{tt}[Renderer=HarfBuzz]{Noto Sans Bengali}
\babelfont[arabic]{rm}[Renderer=HarfBuzz,Scale=1.1]{Noto Naskh Arabic}

% Section styling
\titleformat{\section}{\Large\bfseries\color{emerald}}{\thesection.}{0.5em}{}
\titleformat{\subsection}{\large\bfseries\color{emerald!85!black}}{\thesubsection.}{0.5em}{}
\titleformat{\subsubsection}{\normalsize\bfseries\color{emerald!70!black}}{\thesubsubsection.}{0.5em}{}
\titlespacing*{\section}{0pt}{18pt}{8pt}

% Fancy Header/Footer setup
\pagestyle{fancy}
\fancyhf{} % clear all header and footer fields
\fancyhead[L]{\color{emerald}\small\textbf{\nouppercase{\leftmark}}}
\fancyfoot[L]{\scriptsize\color{black!50}{\lat{AI}}-সংকলিত, আলেম-যাচাইকৃত নয়}
\fancyfoot[C]{\color{emerald!70!black}\thepage}
\renewcommand{\headrulewidth}{0.4pt}
\renewcommand{\headrule}{\hbox to\headwidth{\color{emerald}\leaders\hrule height \headrulewidth\hfill}}
\renewcommand{\footrulewidth}{0pt}

% Custom boxes
% 1. Ayat/Hadith Box
\newtcolorbox{ayatbox}[1][]{
    enhanced, breakable, 
    colback=emerald!5, 
    colframe=emerald, 
    boxrule=0.5pt, 
    arc=3pt, 
    left=10pt, right=10pt, top=10pt, bottom=10pt,
    #1
}

% 2. Quote/Translation Box
\newtcolorbox{quotebox}[1][]{
    enhanced, breakable,
    frame hidden,
    colback=lightgray,
    borderline west={3pt}{0pt}{emerald},
    left=15pt, right=10pt, top=10pt, bottom=10pt,
    sharp corners,
    #1
}

% 3. AI-generated notice box
\newtcolorbox{warnbox}[1][]{
    enhanced, breakable,
    colback=red!4,
    colframe=darkred,
    boxrule=0.8pt,
    arc=3pt,
    left=10pt, right=10pt, top=10pt, bottom=10pt,
    #1
}

% Commands
\newcommand{\arab}[1]{\foreignlanguage{arabic}{#1}}
\newcommand{\kib}[1]{\textcolor{emerald}{\textbf{#1}}}

% Latin (English) fallback: Noto Serif Bengali has NO Latin glyphs
\newfontfamily\latfont[Renderer=HarfBuzz]{Noto Serif}
\newcommand{\lat}[1]{{\latfont #1}}

\setlength{\parskip}{5pt}
\setlength{\parindent}{0pt}

% Fix for missing bullet in Noto Serif Bengali
\renewcommand{\labelitemi}{$\bullet$}



\begin{document}
\begin{center}{\Large\bfseries\color{emerald} 06-দিমাম-নবীজিকে-চিনতে-পারেননি}\end{center}
\vspace{1em}
\section*{ঘটনা ৬: আবু হুরাইরা (রাঃ) — অপরিচিত ব্যক্তি নবীজিকে চিনতে পারেনি}

\begin{warnbox}
 \textbf{গুরুত্বপূর্ণ নোটিশ (\lat{Important} \lat{Notice}):} এই কন্টেন্টটি \lat{AI} (কৃত্রিম বুদ্ধিমত্তা) দিয়ে প্রস্তুত ও সংকলিত। \lat{AI} কঠোর সূত্র-মান নীতি মেনে শুধু নির্ভরযোগ্য উৎস থেকে তথ্য সংগ্রহ করেছে — কেবল সহীহ/হাসান হাদিস ও সহীহ সনদের আসার রাখা হয়েছে, দুর্বল সনদ বাদ দেওয়া হয়েছে। তবে এটি কোনো আলেম বা মুহাদ্দিস কর্তৃক যাচাইকৃত নয়।
\end{warnbox}


\medskip\hrule\medskip

\subsection*{১. সূত্র কার্ড}

\begin{table}[h]\centering\small
\begin{tabular}{ll}
বিষয় & বিবরণ \\
\hline
\textbf{বর্ণনাকারী} & আবু হুরাইরা (রাঃ) থেকে \\
\textbf{গ্রন্থ} & সহীহ বুখারি, হাদিস ৬৩; সুনানে আবু দাউদ, হাদিস ৪৮৬; সুনানে নাসাই ২৬৮৭ \\
\textbf{অধ্যায়} & কিতাবুল ইলম, বাব: কাউকে প্রশ্ন করা যে, তোমাদের মধ্যে অমুক কে? \\
\textbf{সনদের মান} & \textbf{সহীহ} \\
\textbf{মূল বিষয়} & নবীজির বিনয়ী ও সাধারণ জীবন — বাইরের কেউ তাকে আলাদা করে চিনতে পারত না \\
\hline
\end{tabular}
\end{table}

\medskip\hrule\medskip

\subsection*{২. পটভূমি (কে, কখন, কোথায়, কেন)}

এটি মদীনায় ইসলামি \lat{state} (রাষ্ট্র) প্রতিষ্ঠার পরের ঘটনা। নবীজি (সাল্লাল্লাহু আলাইহি ওয়া সাল্লাম) সাধারণত মসজিদে সাহাবাদের সাথে বসে থাকতেন — কোনোরকম রাজকীয় ভাব, আলাদা আসন বা বিশেষ \lat{dignity} (সম্মান)র ব্যবস্থা ছাড়াই।

একদিন \textbf{দিমাম ইবনে সা'লাবা} নামের এক বেদুইন (মরুচারী) — যিনি সা'দ ইবনে বকর \lat{tribe} (গোত্র)র \lat{representative} (প্রতিনিধি) — মদীনায় এলেন। তিনি প্রথমবার নবীজিকে দেখছেন। মসজিদে প্রবেশ করে উট বসিয়ে, বাঁধলেন এবং সাহাবাদের উদ্দেশ্যে বললেন:

\textbf{\lat{“}তোমাদের মধ্যে মুহাম্মদ কে?\lat{”}}

সাহাবারা ইশারা করে বললেন: \textbf{\lat{“}এই যে — সাদা চামড়ার যে লোকটি হেলান দিয়ে বসে আছে।\lat{”}} — এই ছোট্ট একটি দৃশ্য নবীজির \lat{humility} (বিনয়)র সবচেয়ে বড় \lat{evidence} (প্রমাণ)।

\medskip\hrule\medskip

\subsection*{৩. পূর্ণ ঘটনার বর্ণনা}

আবু হুরাইরা (রাঃ) বলেন:

\begin{quote}
\textbf{\lat{“}আমরা নবীজির (সাল্লাল্লাহু আলাইহি ওয়া সাল্লাম) সাথে মসজিদে বসা ছিলাম। তখন এক লোক উটে চড়ে এসে মসজিদে উটটি বসিয়ে দিল, এরপর উটের পা বেঁধে বলল: "তোমাদের মধ্যে মুহাম্মদ কে?"} \\
\textbf{নবীজি তখন আমাদের মাঝে হেলান দিয়ে বসা ছিলেন। আমরা বললাম: "এই — সাদা লোকটি, যে হেলান দিয়ে আছে।"} \\
\textbf{লোকটি বলল: "হে আবদুল মুত্তালিবের ছেলে!"} \\
\textbf{নবীজি বললেন: "আমি তোমার প্রশ্নের জবাব দিচ্ছি।"} \\
\textbf{লোকটি বলল: "আমি তোমাকে কিছু প্রশ্ন করব এবং কঠোরভাবে প্রশ্ন করব — তুমি রাগ করবে না।"} \\
\textbf{নবীজি বললেন: "তুমি যা ইচ্ছা জিজ্ঞেস করো।"} \\
\textbf{লোকটি বলল: "আমি তোমাকে তোমার রবের কসম দিয়ে জিজ্ঞেস করছি — যিনি তোমার পূর্ববর্তীদের রব — আল্লাহ কি তোমাকে সব মানুষের কাছে রাসূল করে পাঠিয়েছেন?" নবীজি বললেন: "আল্লাহর কসম, হ্যাঁ।"} \\
\textbf{...এভাবে সে পাঁচ ওয়াক্ত নামাজ, রমজানের রোজা, যাকাত ইত্যাদির কথা প্রশ্ন করল, নবীজি প্রতিটির জবাবে বললেন: "হ্যাঁ, আল্লাহই আদেশ করেছেন।"} \\
\textbf{সবশেষে লোকটি বলল: "আমি ঈমান আনলাম আপনার আনীত সব কিছুর উপর। আমি আমার গোত্রের পক্ষ থেকে প্রতিনিধি — আমি দিমাম ইবনে সা'লাবা, সা'দ ইবনে বকর গোত্রের ভাই।"\lat{”}}
\end{quote}

\medskip\hrule\medskip

\subsection*{৪. ধাপে ধাপে বিশ্লেষণ}

\textbf{ধাপ ১ — বেদুইনের \lat{arrival} (আগমন):} দিমাম — মক্কা-মদীনা থেকে দূরবর্তী মরু এলাকার মানুষ — নবীজিকে খুঁজতে মদীনায় এলেন।

\textbf{ধাপ ২ — \lat{“}তোমাদের মধ্যে মুহাম্মদ কে?\lat{”}:} সাহাবাদের মধ্যে হেলান দিয়ে বসা নবীজিকে আলাদা করে \lat{recognize} (চেনা) করা গেল না। বাইরের একজন মানুষের কাছে তিনি অন্য সাহাবিদের মতোই \lat{ordinary} (সাধারণ) মানুষ।

\textbf{ধাপ ৩ — সাহাবাদের \lat{gesture} (ইশারা):} সাহাবারা বললেন — \lat{“}সাদা চামড়ার হেলান দেওয়া লোকটি।\lat{”} — অর্থাৎ নবীজিকে কোনো বিশেষ পোশাক, আসন বা \lat{sign} (চিহ্ন) দিয়ে আলাদা করা হতো না।

\textbf{ধাপ ৪ — কঠোর প্রশ্নে নবীজির সহনশীলতা:} দিমাম \lat{harsh} (কঠোর) ভাষায় প্রশ্ন করলেন; নবীজি রাগ করলেন না — বরং \lat{patience} (ধৈর্য)র সাথে প্রতিটি প্রশ্নের জবাব দিলেন।

\textbf{ধাপ ৫ — ইসলামের \lat{core} (মূল) বিষয়:} দিমাম নামাজ, রোজা, যাকাত — ইসলামের \lat{fundamental} (মৌলিক) বিধান সম্পর্কে প্রশ্ন করলেন; জবাব পেয়ে তিনি ইসলাম গ্রহণ করলেন।

\medskip\hrule\medskip

\subsection*{৫. শব্দের ব্যাখ্যা}

\begin{table}[h]\centering\small
\begin{tabular}{ll}
শব্দ & অর্থ \\
\hline
\textbf{আ'রাবী (\arab{أعرابي})} & বেদুইন — মরুভূমির বাসিন্দা, শহরের সভ্যতার বাইরের মানুষ \\
\textbf{মুত্তাকি' (\arab{متكئ})} & হেলান দেওয়া — আরামে বসা \\
\textbf{আবইয়ায (\arab{أبيض})} & সাদা/উজ্জ্বল — এখানে নবীজির উজ্জ্বল ত্বকের বর্ণনা \\
\textbf{আনাখা (\arab{أناخ})} & উট বসিয়ে বিশ্রাম দেওয়া \\
\hline
\end{tabular}
\end{table}

\textbf{এই দৃশ্যের তাৎপর্য:} যে ব্যক্তি সৃষ্টির সেরা মানুষ, তার জীবনে রাজকীয় কোনো আভাস ছিল না। মসজিদে সাহাবাদের সাথে মিশে থাকা, সাধারণ আসন, সাধারণ পোশাক — এই বিনয়ই তার মর্যাদা বাড়িয়েছে।

\medskip\hrule\medskip

\subsection*{৬. সংশ্লিষ্ট হাদিস ও আয়াত}

\textbf{হাদিস (বুখারি ৫৫৬৫):} আনাস (রাঃ) বলেন: \textbf{\lat{“}নবীজি (সা.) বাজারে যেতেন, রোগীদের দেখতে যেতেন, জানাজার পেছনে হাঁটতেন, দাসের দাওয়াত কবুল করতেন।\lat{”}}

\textbf{হাদিস (বুখারি ৬০৭২):} আনাস (রাঃ) বলেন: \textbf{\lat{“}মদীনার কোনো দাসী চাইলে নবীজির হাত ধরে নিজের প্রয়োজনে নিয়ে যেতে পারত।\lat{”}} — নবীজির চরম সহজলভ্যতা ও বিনয়।

\textbf{আয়াত (সূরা আশ-শুআরা ২৬:২১৫):}
\begin{quote}
\textbf{\lat{“}আর যারা আপনার অনুসরণ করেছে, তাদের প্রতি বিনয়ের পাখা নামিয়ে দিন।\lat{”}}
\end{quote}

\textbf{আয়াত (সূরা আত-তাওবা ৯:১২৮):}
\begin{quote}
\textbf{\lat{“}আপনার কাছে এমন একজন রাসূল এসেছেন যিনি তোমাদেরই মধ্য থেকে; তোমাদের কষ্ট তাকে কষ্ট দেয়, তিনি তোমাদের মঙ্গলকামী।\lat{”}}
\end{quote}

\medskip\hrule\medskip

\subsection*{৭. গভীর শিক্ষা}

\textbf{১. প্রকৃত মর্যাদা \lat{humility} (বিনয়)তে।}
নবীজি দুনিয়ার সবচেয়ে \lat{honoured} (মর্যাদাবান) মানুষ — কিন্তু তাঁর মর্যাদা ছিল \lat{humility}তে, রাজকীয় প্রদর্শনে নয়। বাইরের কেউ তাকে সাহাবাদের থেকে আলাদা চিনতে পারত না। এটি শিক্ষা দেয় — বড়ত্বের জন্য \lat{display} (প্রদর্শন)র দরকার নেই।

\textbf{২. নেতার \lat{accessibility} (সহজলভ্যতা)।}
যিনি নিজেকে মানুষের কাছে \lat{accessible} (সহজলভ্য) করেন, মানুষ তাঁকে ভালোবাসে ও বিশ্বাস করে। নবীজি মসজিদে সাহাবাদের সাথে বসতেন, বাজারে যেতেন — মানুষ তাঁকে যে কোনো সময় পেত।

\textbf{৩. \lat{Arrogant} (অহংকারী) নেতা বনাম \lat{Humble} (বিনয়ী) নেতা।}
\lat{Arrogant} (অহংকারী) নেতা নিজেকে আলাদা আসনে বসায়, \lat{respect} (সম্মান)র দাবি করে। \lat{Humble} (বিনয়ী) নেতা মানুষের মাঝে মিশে থাকে — নবীজি ছিলেন এর সর্বোচ্চ \lat{example} (উদাহরণ)।

\textbf{৪. \lat{Knowledge} (জ্ঞান) ও ঈমানের সহজ পথ।}
দিমামের মতো \lat{ordinary} মানুষও সরলভাবে প্রশ্ন করে ইসলাম বুঝতে পেরেছে। \lat{Arrogance} (অহংকার) ও জটিলতা নয়, \lat{simplicity} (সরলতা) ও \lat{humility} জ্ঞানের দরজা খুলে দেয়।

\textbf{৫. সাহাবাদের সাথে নবীজির \lat{relationship} (সম্পর্ক)।}
সাহাবারা নবীজিকে এত কাছে থেকে দেখেছেন, তাদের \lat{testimony} (সাক্ষ্য)ই আমাদের কাছে নবীজির \lat{humility} এর প্রামাণ্য দলিল।

\medskip\hrule\medskip

\subsection*{৮. জীবনে প্রয়োগের পদ্ধতি}

\begin{enumerate}
\item \textbf{নিজেকে সহজলভ্য করুন:} পরিবার, \lat{colleagues} (সহকর্মী) ও \lat{neighbours} (প্রতিবেশী)দের কাছে সর্বদা পৌঁছানোর মতো থাকুন — \lat{arrogance} (অহংকার)র উঁচু প্রাচীর বানাবেন না।
\item \textbf{প্রশ্ন করার \lat{courage} (সাহস):} দিমামের মতো জিজ্ঞেস করতে ভয় পাবেন না; \lat{knowledge} (জ্ঞান) অর্জনে লজ্জা ও \lat{arrogance} বাধা।
\item \textbf{\lat{Display} (প্রদর্শন) নয়, \lat{humility}:} বড় পদ বা মর্যাদায় থাকলে রাজকীয় ভাব নয়, মানুষের সাথে মিশে থাকুন।
\item \textbf{অন্যের সাথে \lat{equal} (সমান) আচরণ:} নবীজির মতো ছোট-বড় সবার দাওয়াত গ্রহণ করুন, দাস-ভৃত্যের সাথে নিজের মতো আচরণ করুন।
\item \textbf{নরম ও সহনশীল জবাব:} কেউ কঠোর প্রশ্ন করলে \lat{anger} (রাগ) নয়, \lat{patience} (ধৈর্য)র সাথে জবাব দিন — নবীজির সুন্নাত।
\end{enumerate}
\medskip\hrule\medskip

\subsection*{৯. রেফারেন্স}

\begin{itemize}
\item সহীহ বুখারি, হাদিস ৬৩ (কিতাবুল ইলম)
\item সুনানে আবু দাউদ, হাদিস ৪৮৬
\item সুনানে নাসাই, হাদিস ২৬৮৭
\item সহীহ বুখারি, হাদিস ৫৫৬৫ ও ৬০৭২ (নবীজির সহজলভ্যতা)
\item সূরা আশ-শুআরা, আয়াত ২১৫; সূরা আত-তাওবা, আয়াত ১২৮
\end{itemize}

\end{document}
